\documentclass[11pt]{article}
\usepackage[margin=0.9in]{geometry}
\usepackage{amsmath,amssymb,amsthm,booktabs}
\usepackage{mathrsfs}
\usepackage[T1]{fontenc}
\usepackage{lmodern,microtype}
\usepackage[colorlinks=true,linkcolor=blue,urlcolor=blue]{hyperref}
\newtheorem{theorem}{Theorem}[section]
\newtheorem{lemma}[theorem]{Lemma}
\newtheorem{proposition}[theorem]{Proposition}
\newtheorem{corollary}[theorem]{Corollary}
\theoremstyle{remark}
\newtheorem{remark}[theorem]{Remark}
\newcommand{\norm}[1]{\left\|#1\right\|}
\title{Toward a four-dimensional Yang--Mills construction\\
\large Notebook 16: local nonlinear responses and the classical coarse action}
\author{Research working document}
\date{16 October 2026}
\begin{document}
\maketitle
\begin{abstract}
Starting from the finite-regulator small-field minimizer of
Notebook 15, we prove exponential locality of its response to
coarse data at small nonconstant backgrounds. The bounds are
uniform in refinement and periodic volume, retain the same
physical H\"older exponent, and allow one weighted direction
at every analytic derivative order. The extra input is a local
physical-transpose derivative estimate for the nonlinear
constraint. An exact envelope identity then identifies the
derivative of the minimized classical action with the coarse
multiplier. This supplies a local nonlinear coarse Hessian,
its small perturbation from the flat Maxwell precision, and a
classical gauge identity. A separate exact Gaussian trace
calculation shows that a fixed physical smooth ball has vanishing
probability when the coupling decreases more slowly than the
square of the mesh spacing. This identifies a concrete need for
rescaled fluctuation domains. Notebook 17 supplies the resulting
fixed-step Gaussian domain and joint tails; Notebook 18 supplies
an overlapping-interaction integration theorem under explicit local
coefficient hypotheses. The complete model-specific integration
and iteration remain open.
\end{abstract}

\noindent\textbf{Status.}
Internally audited research proofs, with no external certification
or novelty claim. The prerequisite finite-regulator construction
and its ordinary-constraint transfer are given in Notebook 15.
The source provenance remains the deposited transcriptions at
\href{https://github.com/ethanwsussman/Yang--Mills-LLM-experiment/tree/e392ef45013a581664bab76b90867170b30fe10e}{commit
\texttt{e392ef45013a581664bab76b90867170b30fe10e}}.
No new original-paper verification is asserted here.

\section{The precise target}
\label{nr16-target}

Work with $SU(N_c)$ on compatible finite periodic lattices,
fine spacing $\eta=L^{-j}$ for fixed odd $L\ge3$, coarse
spacing one, and fixed $0<\alpha<1$. The background coarse
potential $B$ is arbitrary in a sufficiently small global
supremum ball; it need not be constant or decay at infinity.
The selected fine potential $A_*(B)$ uses the corrected
constraint and fixed flat gauge of Notebook 15. All derivatives
of this chosen representative below refer to that selection.

The ordinary-feasible representative from Notebook 15 is gauge
equivalent and has a separately controlled response through its
explicit gauge formula. Its leading extension generally differs
from $H$, while its minimized action value is exactly the same.
Keeping these coordinates distinct gives local response theorems
for both specified representatives and the ordinary-gauge
interpretation of the classical action.

The main analytic estimate puts an exponential weight on one
coarse direction while keeping the other directions and the
background unweighted. It yields decay away from a localized
insertion and uniform higher derivatives. A fixed-exponent
connected-tree bound involving all insertion sites at once is
a stronger statement and is not inferred from this estimate.


% No preamble. Notebook 16 candidate module. Unique labels: wr-.
% Requires amsmath, amssymb, amsthm and a one-argument \norm macro.

\section{Exponential locality of nonlinear minimizer responses}
\label{wr-section}

Fix $0<\alpha<1$, an odd blocking ratio $L\ge3$, refinement
$M=L^j$, fine spacing $\eta=M^{-1}$, and final spacing one.
The primary statements are on finite compatible periodic
regulators, with constants independent of refinement and volume.
The coarse background may be nonconstant, but it lies in a fixed
small complex bounded-field ball. We study the corrected-constraint
minimizer $A(B)$ in its fixed flat gauge. A gauge-equivalent
representative chosen by a different construction does not acquire
these locality estimates merely from gauge equivalence.

\subsection{Weighted spaces and the fixed flat inputs}
\label{wr-spaces}

Write
\[
 \mathcal F(A)=QA+\mathcal R(A),\qquad
 D_A=D\mathcal R(A),\qquad
 \mathcal X=\ker Q\cap\ker(RD^T).
\]
Here $R$ is the fixed scalar gauge projection and
$\mathcal R$ is the nonlinear averaging remainder. All transposes
use physical measure. Let $\mathcal N_\alpha$ be the complete
potential norm consisting of the supremum, the suprema of all
physical first differences, and their local $\alpha$-H\"older
seminorms. Fix unit-cell collars of bounded physical size.

Choose a sufficiently small fixed $\kappa>0$ allowed by the flat
kernel and constrained Schauder estimates. An admissible positive
cell weight satisfies
\begin{equation}\label{wr-weight}
 \frac{w(n)}{w(m)}\le e^{\kappa d(n,m)}.
\end{equation}
Weighted norms take the supremum of $w(n)$ times the appropriate
local norm on the collar of cell $n$. Fine derivatives are taken
on complete fields before these norms are evaluated. No sharp cell
cutoff or derivative of the weight is inserted. Periodic distance
is used on a torus, and fixed collar enlargements only change
constants.

On the infinite lattice a decaying admissible weight can allow
unbounded weighted directions. Literal Fr\'echet responses of the
original bounded-field map are first taken on globally bounded
directions with finite weighted norm. The linearized formulas
below then define bounded extensions to the full weighted spaces.
Higher-response symmetry is initially understood on bounded
directions. The localized corollaries use growing weights
$w_S\ge1$ and require no such extension.

The proved flat operator estimates, uniform for all these weights,
are
\begin{align}
 H &: \ell^\infty_{c,w}\longrightarrow\mathcal N_{\alpha,w},&
 \mathcal C &: \ell^\infty_{f,w}\longrightarrow\mathcal N_{\alpha,w},
                                                       \label{wr-flat-a}\\
 \mathcal C S_{\rm jet}^T &: C^\alpha_{f,w}\longrightarrow\mathcal N_{\alpha,w},
 &q &: \ell^\infty_{c,w}\longrightarrow\ell^\infty_{c,w},
                                                       \label{wr-flat-b}\\
 H^T &: \ell^\infty_{f,w}\longrightarrow\ell^\infty_{c,w},&
 (S_{\rm jet}H)^T &: \ell^\infty_{f,w}\longrightarrow\ell^\infty_{c,w}.
                                                       \label{wr-flat-c}
\end{align}
The symbol $c$ denotes coarse fields and $f$ fine fields.
The exponent in \eqref{wr-flat-b} is unchanged; the separately
proved constrained biharmonic cancellation is needed for this.
The transpose bounds in \eqref{wr-flat-c} use the physical volume
one of a fine cell and summable derivative-extension kernels.

Write the nonlinear Wilson force as
$N(A)=S_{\rm jet}^Tf_\eta(S_{\rm jet}A)$. The local analytic
coefficient and the complete-field H\"older product estimate give
\begin{equation}\label{wr-force}
 \norm{D[f_\eta(S_{\rm jet}A)]V}_{C^\alpha_w}
 \le C\norm{A}_{\mathcal N_\alpha}
                   \norm{V}_{\mathcal N_{\alpha,w}}.
\end{equation}
Thus $\mathcal C N'(A)$ and $H^TN'(A)$ have the corresponding
weighted norms bounded by $C_\kappa\norm{A}_{\mathcal N_\alpha}$.
The second operator is evaluated using $(S_{\rm jet}H)^T$;
no rough nonlinear coefficient is differentiated.

\subsection{A local estimate for the differentiated transpose}
\label{wr-local-transpose}

The full corrected average has a physical dependency radius
$R_*$ independent of $j$. One may use the deposited two-block
property, or sum the one-step radii using
$\sum_{t<j}h_t\le(L-1)^{-1}$. Increase $R_*$ by a fixed cell
collar margin. The established absolute physical row estimates are
\begin{equation}\label{wr-row-estimates}
 \norm{D_A}_{\infty\to\infty}\le c\norm{A}_\infty,
 \qquad
 \norm{D_A^T}_{\infty\to\infty}\le c_*\norm{A}_\infty.
\end{equation}
These inequalities hold in the absolute row norm as well.
For the second inequality the one-step physical column factor
is $L^d$, giving
$N_LK_0L^{d-1}h\norm{A}_\infty$ before summing the spacings.
It is not a comparison with the number of fine variables.

Finite dependency immediately yields
\begin{equation}\label{wr-weighted-rows}
 \begin{split}
 \norm{D_AV}_{\infty,w}
 &\le ce^{\kappa R_*}\norm{A}_\infty\norm{V}_{\infty,w},\\
 \norm{D_A^Tm}_{\infty,w}
 &\le c_*e^{\kappa R_*}\norm{A}_\infty\norm{m}_{\infty,w}.
 \end{split}
\end{equation}
The key additional estimate places the weight on a background
variation while leaving the coarse multiplier unweighted:
\begin{equation}\label{wr-mixed-transpose}
 \boxed{\quad
 \norm{(D_A^T)'[V]m}_{\infty,w}
 \le C_*e^{2\kappa R_*}\norm{V}_{\infty,w}\norm{m}_\infty.
 \quad}
\end{equation}
The constant is uniform on a fixed smaller complex background
ball.

\begin{proof}[Proof of \eqref{wr-mixed-transpose}]
Fix an output fine bond $b$. Its row in $D_A^T$ reads only coarse
bonds within $R_*$ of $b$. Each associated coefficient depends on
$A$ only in the dependency set of that coarse bond, hence within
$2R_*$ of $b$. The map from the bounded background ball into this
physical absolute row space is analytic and uniformly bounded
by \eqref{wr-row-estimates}. Cauchy on a fixed larger ball bounds
its derivative by $C_*\norm{V}_\infty$.

For this one row restrict $V$ to the union of those dependency
sets. The restriction does not change the row derivative, and its
supremum is at most
$e^{2\kappa R_*}w(b)^{-1}\norm{V}_{\infty,w}$, with the fixed
collar convention absorbed in $R_*$. This restriction is used
only in a bounded-value coefficient map, where no field derivative
is taken. Sum the resulting physical row against bounded $m$,
multiply by $w(b)$, and take the supremum over $b$.
The same argument works with lifted multiplicities on small tori.
\end{proof}

The doubled dependency radius in this proof is essential.
Cauchy is applied to the uniformly bounded physical row, not
separately to every fine coefficient followed by an uncontrolled
count. Higher coefficient derivatives admit the analogous estimate
with one weighted direction and the remaining directions bounded,
with factorial constants from the analytic margin. The earlier
two-Hilbert-direction Hessian incidence estimate is consistent
with this bound but does not replace its proof.

\subsection{Weighted variation of the chart and multiplier}
\label{wr-variation}

The exact chart and multiplier are
\begin{equation}\label{wr-chart-multiplier}
 \begin{split}
 A&=\xi+H[B-\mathcal R(A)],\\
 m&=(I+H^TD_A^T)^{-1}
             \{q[B-\mathcal R(A)]+H^TN(A)\}.
 \end{split}
\end{equation}
Let $r$ bound $\norm{A}_{\mathcal N_\alpha}$ and $b$ bound
$\norm{B}_\infty$. The established unweighted estimate is
\begin{equation}\label{wr-background-multiplier}
 \norm{m}_\infty\le C(b+r^2).
\end{equation}
Neither $A$ nor $m$ is required to have a finite growing-weight
norm. Only their variations will carry a weight.

Shrink the common background radius once more, depending on
$\kappa$, so that uniformly in admissible $w$,
\begin{equation}\label{wr-inverse-smallness}
 \norm{HD_A}_{\mathcal N_{\alpha,w}\to\mathcal N_{\alpha,w}}<\frac12,
 \qquad
 \norm{H^TD_A^T}_{\infty,w\to\infty,w}<\frac12.
\end{equation}
These inequalities follow from \eqref{wr-flat-a},
\eqref{wr-flat-c}, and \eqref{wr-weighted-rows}. Their radii
are independent of source location, refinement, and volume.

For $z=\delta\xi$, $b_1=\delta B$, and $v=\delta A$, exact
differentiation gives
\begin{equation}\label{wr-chart-variation}
 v=P_A(z+Hb_1),\qquad P_A=(I+HD_A)^{-1},\qquad
 \norm{v}_{\mathcal N_{\alpha,w}}
 \le2\bigl(\norm{z}_{\mathcal N_{\alpha,w}}
                       +h_\kappa\norm{b_1}_{\infty,w}\bigr).
\end{equation}
Differentiate the defining equation for $m$ in
\eqref{wr-chart-multiplier}. This preserves the noncommuting
order and gives
\begin{equation}\label{wr-multiplier-variation}
 \boxed{\quad
 \delta m=(I+H^TD_A^T)^{-1}
 \{q(b_1-D_Av)+H^TN'(A)v-H^T(D_A^T)'[v]m\}.
 \quad}
\end{equation}
Using \eqref{wr-force}, \eqref{wr-weighted-rows},
\eqref{wr-mixed-transpose}, and
\eqref{wr-background-multiplier},
\begin{equation}\label{wr-multiplier-variation-bound}
 \norm{\delta m}_{\infty,w}
 \le C_\kappa\bigl(\norm{b_1}_{\infty,w}
        +(r+b+r^2)\norm{v}_{\mathcal N_{\alpha,w}}\bigr).
\end{equation}
In particular the last term of
\eqref{wr-multiplier-variation} needs the mixed transpose
estimate; no weighted norm of the global background multiplier
has been assumed.

\subsection{The weighted response inverse}
\label{wr-response-inverse}

The selected minimizer solves
\begin{equation}\label{wr-fixed-point}
 \xi=\Phi_B(\xi)=-\mathcal C N(A)+\mathcal C D_A^Tm.
\end{equation}
Its variation is exactly
\begin{equation}\label{wr-phi-variation}
 \delta\Phi=-\mathcal C N'(A)v
       +\mathcal C\{(D_A^T)'[v]m+D_A^T\delta m\}.
\end{equation}
The preceding estimates imply, for $r<1$,
\begin{equation}\label{wr-phi-bound}
 \norm{\delta\Phi}_{\mathcal N_{\alpha,w}}
 \le C_\kappa\bigl((r+b)\norm{v}_{\mathcal N_{\alpha,w}}
                            +r\norm{b_1}_{\infty,w}\bigr).
\end{equation}
For $b_1=0$, combine this with
\eqref{wr-chart-variation}. Choose the chart and coarse
background radii so that $2C_\kappa(r+b)<1/2$. Then
\begin{equation}\label{wr-weighted-contraction}
 \norm{D_\xi\Phi_B}_{\mathcal N_{\alpha,w}\cap\mathcal X
                         \to\mathcal N_{\alpha,w}\cap\mathcal X}
 <\frac12.
\end{equation}
Every term lies in $\mathcal X$ because it contains $\mathcal C$.
The constrained space is closed on each finite regulator. The
estimate does not use a bounded-field norm for the separate gauge
projection.

The derivative of the minimizer is therefore obtained by the
weighted Neumann inverse
\begin{equation}\label{wr-derivative-formula}
 z=(I-D_\xi\Phi_B)^{-1}D_B\Phi_B[b_1],\qquad
 D_BA(B)[b_1]=P_A(z+Hb_1).
\end{equation}
At the solution $\norm{A(B)}_{\mathcal N_\alpha}
\le C\norm{B}_\infty$. Equations \eqref{wr-phi-bound} and
\eqref{wr-derivative-formula} show that $z$ is
$O(\norm{B}_\infty\norm{b_1}_{\infty,w})$.
Also $P_A-I=-P_AHD_A$ is $O(\norm{B}_\infty)$ in the
weighted norm. Consequently
\begin{equation}\label{wr-main-response}
 \boxed{\begin{aligned}
 \norm{D_BA(B)[b_1]}_{\mathcal N_{\alpha,w}}
 &\le C_\kappa\norm{b_1}_{\infty,w},\\
 \norm{(D_BA(B)-H)[b_1]}_{\mathcal N_{\alpha,w}}
 &\le C_\kappa\norm{B}_\infty\norm{b_1}_{\infty,w}.
 \end{aligned}}
\end{equation}
This is a response estimate at arbitrary small nonconstant
backgrounds derived using fixed flat operators and nonlinear
coefficients. No nonflat source Schauder theorem is assumed.

\subsection{Coarse multiplier responses}
\label{wr-multiplier-responses}

Now set $m(B)=m(A(B),B)$ and
$v=D_BA(B)[b_1]$ in \eqref{wr-multiplier-variation}. It immediately
gives
\begin{equation}\label{wr-dm-bound}
 \norm{D_Bm(B)[b_1]}_{\infty,w}
 \le C_\kappa\norm{b_1}_{\infty,w}.
\end{equation}
More explicitly, putting $M_A=I+H^TD_A^T$ and subtracting
$qb_1$ gives the exact expression
\begin{equation}\label{wr-dm-remainder}
 \begin{split}
 (D_Bm(B)-q)b_1={}&(M_A^{-1}-I)qb_1\\
 &+M_A^{-1}\{-qD_Av+H^TN'(A)v-H^T(D_A^T)'[v]m\}.
 \end{split}
\end{equation}
At the solution the inverse difference, $D_A$, and $N'(A)$
are $O(\norm{B}_\infty)$ in their respective weighted spaces.
Although $(D_A^T)'$ need not vanish at zero, its multiplier
has norm $\norm{m}_\infty=O(\norm{B}_\infty)$, so
\eqref{wr-mixed-transpose} controls that term too. Thus
\begin{equation}\label{wr-dm-near-q}
 \boxed{\quad
 \norm{(D_Bm(B)-q)b_1}_{\infty,w}
 \le C_\kappa\norm{B}_\infty\norm{b_1}_{\infty,w}.
 \quad}
\end{equation}
The multiplier is holomorphic as a global bounded-field map,
and its derivative is holomorphic in these weighted operator
spaces by the displayed formulas and the uniformly convergent
weighted Neumann series.

\subsection{Localized changes and all-order bounds with one weight}
\label{wr-localization}

If $b_1$ is supported in coarse cells $S$, choose
$w_S(n)=e^{\kappa d(n,S)}$. This is admissible. Equation
\eqref{wr-main-response} gives
\begin{equation}\label{wr-local-response}
 \norm{D_BA(B)[b_1]}_{\mathcal N_\alpha(C_n^+)}
 \le C_\kappa e^{-\kappa d(n,S)}\norm{b_1}_\infty.
\end{equation}
Equations \eqref{wr-dm-bound} and \eqref{wr-dm-near-q}
give the analogous coarse multiplier estimates. For two small
backgrounds whose connecting coarse line segment stays inside
the common ball, integrate these derivative bounds to obtain
\begin{equation}\label{wr-finite-difference}
 \norm{A(B')-A(B)}_{\mathcal N_{\alpha,w}}
 +\norm{m(B')-m(B)}_{\infty,w}
 \le C_\kappa\norm{B'-B}_{\infty,w}.
\end{equation}
The individual backgrounds may be globally nondecaying. On the
full lattice truncated admissible weights and passage to their
limit justify the same statement without assuming the difference
in advance belongs to the growing-weight space.

On a smaller common complex background ball,
$B\mapsto D_BA(B)$ and $B\mapsto D_Bm(B)$ are holomorphic
into the weighted operator spaces. On finite regulators this is
finite-dimensional holomorphy with the uniform bounds just
proved. Equivalently it follows from the local coefficient
analyticity, the H\"older product estimates, and the weighted
inverse formulas, which also apply to complete infinite fields.
Background directions in this assertion are unweighted bounded
coarse fields. If the singled-out weighted direction is unbounded
on the infinite lattice, the operator-valued formulas specify the
linearized extension just described, rather than a Fr\'echet
derivative in an inadmissible background direction.

Let $\delta>0$ be the margin to a larger complex background ball
where the uniform weighted derivative bounds hold. Operator-valued
multivariable Cauchy gives, for every $n\ge1$,
\begin{equation}\label{wr-all-orders}
 \boxed{\begin{split}
 &\norm{D_B^n A(B)[b_1,\ldots,b_n]}_{\mathcal N_{\alpha,w}}
       +\norm{D_B^n m(B)[b_1,\ldots,b_n]}_{\infty,w}\\
 &\hspace{5mm}\le C_\kappa(n-1)!
       \left(\frac e\delta\right)^{n-1}
        \norm{b_1}_{\infty,w}\prod_{i=2}^n\norm{b_i}_\infty.
 \end{split}}
\end{equation}
For $n=1$ this is the proved first-response estimate. For $n>1$
use polyradii approaching $\delta/(n-1)$ and
$(n-1)^{n-1}\le e^{n-1}(n-1)!$. Symmetry permits any single
direction to carry the weight.

If all directions are supported in sets $S_i$, applying
\eqref{wr-all-orders} separately with each direction weighted
gives decay $\exp[-\kappa\max_i d(n,S_i)]$ at an output cell
$n$. This is a precise one-weight conclusion. It is not a
fixed-exponent tree-length or ordered-path estimate among all
insertion sites. Such stronger multivertex bounds require an
additional branched expansion or multiweight argument.

\subsection{Scope of the response theorem}
\label{wr-scope}

The exponent $\kappa$ is fixed once, and the allowed background
ball is reduced by a fixed amount to satisfy the weighted
inverse inequalities. Its radius does not depend on refinement,
volume, the insertion location, or the particular admissible
weight. The theorem does not assert that the largest unweighted
existence radius already has this property.

These results concern the selected corrected-constraint minimizer
in the fixed flat gauge and its coarse multiplier. They do not
identify its derivatives with a moving-gauge source minimizer,
establish mixed-history estimates, or prove an interacting
continuum limit. All action-minimum statements remain finite
regulator statements; bounded infinite fields need not have finite
total action.

% No preamble. Unique prefix or-. Uses \norm and standard theorem environments.
% Locality of the explicit ordinary representative, independently audited.

\section{Locality of the chosen ordinary-constraint representative}
\label{or-section}

Gauge equivalence by itself does not preserve locality of a chosen
potential representative. For the explicit gauge used in the ordinary
constraint transfer, however, the required locality can be proved.
Write $A_*(B)$ for the corrected fixed-slice minimizer and retain the
cumulative endpoint gauge $b_j$ and the fixed bump extension
$E_{\rm pt}$ from that transfer. Set
\begin{equation}\label{or-transfer}
 \omega(A)=\frac1i\log b_j(e^{i\eta A})^{-1},\qquad
 \lambda(A)=E_{\rm pt}\omega(A),\qquad
 T(A)=\mathscr A_\eta(A,\lambda(A)),\qquad
 A_{\rm ord}(B)=T(A_*(B)).
\end{equation}
All bases are small in the unweighted global norms. Weights are
placed on variations only. Fix $0<\alpha<1$ and one sufficiently
small positive spatial exponent $\kappa$.

\subsection{The derivative of the gauge transfer is weighted bounded}

The cumulative gauge is uniformly holomorphic in the global
potential supremum norm, has $b_j-I=O(\norm{A}_\infty)$,
and has a fixed final physical dependency radius $R_b$.
The latter follows from the endpoint product and the summable
physical reading radii of the corrected trajectory.
The same properties hold for $\omega$, with $\omega(0)=0$.
For an admissible weight $w(n)/w(m)\le e^{\kappa d(n,m)}$,
local-row Cauchy gives
\begin{equation}\label{or-omega-derivative}
 \norm{D\omega(A)V}_{\infty,w}
      \le C e^{\kappa R_b}\norm{V}_{\infty,w}.
\end{equation}
Indeed, for one coarse output row, restrict the direction to that
row's dependency neighborhood. Its unweighted supremum is bounded
by the right side divided by the output weight. No derivative is
taken of this restricted direction: $\omega$ consumes bounded
potential values. Cauchy is applied to the whole bounded row,
not coefficientwise followed by a fine-dimensional count.
The same argument proves operator-valued holomorphy in the
weighted norm on a smaller unweighted background ball.

The fixed bump interpolation satisfies
\begin{equation}\label{or-point-weighted}
 \norm{E_{\rm pt}c}_{C^{2,\alpha}_w}
          \le C_\kappa\norm{c}_{\infty,w}.
\end{equation}
Each collar reads a fixed number of nearby coarse coefficients.
Bounded smooth bump derivatives through order three, weight
comparison on that neighborhood, and sampling give this estimate,
uniformly in refinement and compatible periodic volume.
Consequently
\begin{equation}\label{or-lambda}
 \norm{D\lambda(A)V}_{C^{2,\alpha}_w}
       \le C_\kappa\norm{V}_{\infty,w},\qquad
 \norm{\lambda(A)}_{C^{2,\alpha}}
       \le C\norm{A}_\infty.
\end{equation}
The base estimate in the second inequality is unweighted.

For completeness, the fine gauge action is
\[
 \mathscr A_\eta(A,\lambda)(e)=\frac1{i\eta}\log
 (e^{i\lambda(e_-)}e^{i\eta A(e)}e^{-i\lambda(e_+)}).
\]
Writing $\ell=\lambda(e_-)$ and
$v=(\lambda(e_+)-\lambda(e_-))/\eta$, its group argument is
$e^{i\eta\operatorname{Ad}_{e^{i\ell}}A(e)}
e^{i\ell}e^{-i(\ell+\eta v)}$.
It equals identity at $\eta=0$, so its logarithm divided by
$\eta$ has a uniform analytic jet extension there.
The derivative is therefore a finite sum of products of analytic
coefficients of the base jet $(A,\lambda,D\lambda)$ and the
variation jet $(V,\mu,D\mu)$. The base coefficients have bounded
unweighted $C^{1,\alpha}$ norms. The discrete product estimate
with one weighted factor gives
\begin{equation}\label{or-gauge-derivative}
 \norm{D\mathscr A_\eta(A,\lambda)[V,\mu]}_{N_{\alpha,w}}
 \le C_\kappa\bigl(\norm{V}_{N_{\alpha,w}}
                         +\norm{\mu}_{C^{2,\alpha}_w}\bigr).
\end{equation}
A first output difference uses at most two scalar differences.
The same jet estimate controls its H\"older member; there is no
inverse-mesh loss or weighted-smallness assumption on the base.
Combining \eqref{or-lambda} and \eqref{or-gauge-derivative} proves
\begin{equation}\label{or-DT}
 \norm{DT(A)V}_{N_{\alpha,w}}
                  \le C_\kappa\norm{V}_{N_{\alpha,w}},
\end{equation}
with holomorphic dependence in this weighted operator norm.

\subsection{Responses and the new leading extension}

The corrected minimizer's weighted response theorem and the chain
rule now give
\begin{equation}\label{or-first}
 D_BA_{\rm ord}(B)=DT(A_*(B))D_BA_*(B),\qquad
 \norm{D_BA_{\rm ord}(B)b}_{N_{\alpha,w}}
                 \le C_\kappa\norm{b}_{\infty,w}.
\end{equation}
Its leading extension is
\begin{equation}\label{or-leading}
 H_{\rm ord}=DT(0)H,
\end{equation}
and can differ from $H$. To check the sign explicitly, let
$\beta(A)=\log b_j(e^{i\eta A})/i$ and
$B_{\rm cum}=D\beta(0)$. On the common branch $\omega=-\beta$,
and the gauge action has derivative $V-D\mu$ at zero. Hence
\begin{equation}\label{or-leading-explicit}
 DT(0)=I+DE_{\rm pt}B_{\rm cum},\qquad
 H_{\rm ord}=H+DE_{\rm pt}B_{\rm cum}H.
\end{equation}
Thus $dH_{\rm ord}=dH$.
The exact transfer identity
$\mathcal F_{\rm ord}(T(A))=\mathcal F_{\rm corr}(A)$
also gives $Q_{\rm ord}DT(0)=Q$ and
$Q_{\rm ord}H_{\rm ord}=I$, where $Q_{\rm ord}$ is the
derivative of the ordinary homogeneous constraint at zero.
Here $\mathcal F_{\rm ord}(A)=\log M^{(j)}(e^{i\eta A})/i$
and $\mathcal F_{\rm corr}=\mathcal F$; the complete norm
$N_\alpha$ is the same as $\mathcal N_\alpha$ used above.

The response family in \eqref{or-first} is holomorphic on a
common unweighted complex coarse ball with a uniform weighted
operator norm. Operator Cauchy on a smaller ball yields
\begin{equation}\label{or-linear-correction}
 \norm{(D_BA_{\rm ord}(B)-H_{\rm ord})b}_{N_{\alpha,w}}
       \le C_\kappa\norm{B}_\infty\norm{b}_{\infty,w}.
\end{equation}
If $\delta$ is a fixed analytic margin, then for every $n\ge1$,
\begin{equation}\label{or-higher}
 \norm{D_B^nA_{\rm ord}(B)[b_1,\ldots,b_n]}_{N_{\alpha,w}}
 \le C_\kappa(n-1)!\left(\frac e\delta\right)^{n-1}
       \norm{b_1}_{\infty,w}\prod_{i=2}^n\norm{b_i}_\infty.
\end{equation}
For $n>1$ apply multivariable Cauchy to the first-response
operator with $n-1$ unweighted directions and use
$(n-1)^{n-1}\le e^{n-1}(n-1)!$. The case $n=1$ is
\eqref{or-first}. Symmetry permits any selected bounded direction
to carry the weight.

For a direction supported in S, choose
$w_S(n)=e^{\kappa d(n,S)}$. Equations \eqref{or-first} and
\eqref{or-higher} give exponential decay away from S. If all
directions are supported in sets $S_i$, applying the bound
separately to each weighted direction gives the factor
$e^{-\kappa\max_i d(n,S_i)}$. No tree-length estimate is inferred.
Integration along a straight coarse segment gives
\begin{equation}\label{or-finite-change}
 \norm{A_{\rm ord}(B')-A_{\rm ord}(B)}_{N_{\alpha,w}}
       \le C_\kappa\norm{B'-B}_{\infty,w}.
\end{equation}
Thus finite changes supported in S have the same spatial decay.
For complete infinite fields use bounded truncated weights and
pass to the limit. For weights that decay to zero, literal
Fr\'echet responses are first read on bounded directions with
finite weighted norm; the formulas define their bounded extension
to the full weighted space. The support weights satisfy $w_S\ge1$
and require no such qualification.

This proves locality of the particular ordinary representative
defined in \eqref{or-transfer}; the field-dependent gauge has been
estimated rather than treated as harmless by gauge equivalence.
It does not assert locality of every gauge-equivalent representative,
identify the separately moving source gauge, or cover mixed
retained-history geometries. The action-minimum statement remains
the finite homogeneous small-global-field theorem proved earlier.


% No preamble. Notebook 16 classical coarse-action module. Unique labels: ca-.
% Place after the wr- response module. Standard theorem environments and \norm.

\section{The classical coarse action and its gauge identity}
\label{ca-section}

We remain on finite compatible periodic regulators, with fine
spacing $\eta=L^{-j}$ and final coarse spacing one. The action
$\mathcal S_\eta(U)$ is the coupling-free Wilson action in the
normalization of Notebook 15. Set
$\mathcal W_\eta(A)=\mathcal S_\eta(e^{i\eta A})$.
All coarse pairings below use the normalized color trace and
the physical coarse measure, which is one at this scale.
Transposes are bilinear after complexification. The constants
are independent of refinement and periodic volume.

Let $A_*(B)$ be the corrected-constraint minimizer in the fixed
flat gauge slice, and put $U_*(B)=e^{i\eta A_*(B)}$. Its
ordinary-feasible representative $U_{\rm ord}(B)$ is related
to $U_*(B)$ by the genuine gauge transformation constructed in
Notebook 15. Define
\begin{equation}\label{ca-definition}
 V(B)=\mathcal W_\eta(A_*(B))
     =\mathcal S_\eta(U_*(B))
     =\mathcal S_\eta(U_{\rm ord}(B)).
\end{equation}
These are equal holomorphic functions on a uniform complex
coarse supremum ball. Indeed the analytic Wilson density,
written with $P+P^{-1}$, is invariant under complex matrix
gauge conjugation as well. Only the real fields have the
minimum interpretation. The scalar total action scales with
volume; the uniform estimates below state that dependence.

\subsection{The multiplier is the exact derivative of the action}
\label{ca-envelope-section}

Write $\mathcal F(A)=QA+\mathcal R(A)$ and use the exact
chart and flat-slice space
\[
 A=\xi+H[B-\mathcal R(A)],\qquad
 \xi\in\mathcal X=\ker Q\cap\ker(R_0D^T).
\]
Here $R_0$ is the fixed scalar gauge projection, distinct from
the nonlinear remainder $\mathcal R$ and denoted $R$ in the
response module. Put
\begin{equation}\label{ca-chart-differential}
 T_A=(I+HD\mathcal R(A))^{-1},\qquad H_A=T_AH.
\end{equation}
Differentiating the analytic minimizer in a coarse direction $c$
gives the exact formula
\begin{equation}\label{ca-response-decomposition}
 D_BA_*(B)[c]=H_Ac+T_AD_B\xi(B)[c],\qquad A=A_*(B).
\end{equation}
The second summand is a slice tangent because
$D_B\xi(B)[c]\in\mathcal X$.

\begin{proposition}[Classical envelope identity]\label{ca-envelope}
Let $E(A)=D\mathcal W_\eta(A)$ denote the Euler field under
the physical fine pairing. The multiplier from the exact
stationarity construction satisfies
\begin{equation}\label{ca-first-derivative}
 DV(B)[c]=\langle m(B),c\rangle_1,\qquad
 m(B)=H_A^TE(A_*(B)).
\end{equation}
Equivalently,
\begin{equation}\label{ca-multiplier}
 \begin{split}
 m(B)={}&(I+H^TD\mathcal R(A_*)^T)^{-1}\\
 &\quad\cdot\{q[B-\mathcal R(A_*)]+H^TN(A_*)\},
 \end{split}
\end{equation}
where $q=q_M$ is the physical flat Maxwell precision and
$E(A)=KA+N(A)$ is the ordinary-coordinate force.
\end{proposition}

\begin{proof}
Pair \eqref{ca-response-decomposition} with $E(A_*)$.
Slice stationarity makes
$\langle E(A_*),T_AD_B\xi(B)[c]\rangle_\eta=0$.
The remaining term is exactly \eqref{ca-first-derivative}.
This argument needs only slice stationarity, not a separate
full-constraint multiplier theorem or gauge coverage.

The transpose push-through identity gives
$H_A^T=(I+H^TD\mathcal R(A_*)^T)^{-1}H^T$.
The flat identities $H^TK\xi=0$ and $H^TKH=q$ imply
$H^TKA_*=q[B-\mathcal R(A_*)]$, proving
\eqref{ca-multiplier}. The force term is evaluated as
$H^TN(A_*)=(S_{\rm jet}H)^Tf_\eta(S_{\rm jet}A_*)$;
no rough force coefficient is differentiated. All identities
extend bilinearly to the complex ball.
\end{proof}

Put $b=\norm{B}_\infty$. The uniform estimates
$A_*=HB+O(b^2)$, $\mathcal R(A_*)=O(b^2)$,
$H^TN(A_*)=O(b^2)$ and
$\norm{D\mathcal R(A_*)^T}_{\infty\to\infty}=O(b)$
in \eqref{ca-multiplier} prove
\begin{equation}\label{ca-m-near-q}
 \norm{m(B)-qB}_\infty\le Cb^2,
 \qquad m(0)=0,\quad Dm(0)=q.
\end{equation}
The multiplier is holomorphic on a uniform coarse ball.
In particular $D^2V(0)=q$. Its unweighted estimate
$\norm{Dm(B)-q}_{\infty\to\infty}\le Cb$ also follows
directly by Cauchy: for $b>0$ use a direction circle of radius
$b$ inside a smaller common ball, and apply
\eqref{ca-m-near-q} at points of norm at most $2b$.
At $b=0$ use the displayed Taylor coefficient.

Let $N_b$ denote the number of independent positive coarse
bonds. Since $V(0)=0$, radial integration gives
\begin{equation}\label{ca-radial}
 V(B)=\frac12\langle B,qB\rangle_1
       +\int_0^1\langle m(tB)-tqB,B\rangle_1\,dt.
\end{equation}
The fixed normalized color trace bound and
\eqref{ca-m-near-q} therefore yield
\begin{equation}\label{ca-cubic-remainder}
 \boxed{\quad
 \left|V(B)-\frac12\langle B,qB\rangle_1\right|
          \le C N_b\norm{B}_\infty^3.
 \quad}
\end{equation}
This holds on the complex ball. It is a volume-uniform
density estimate, not a volume-independent total-action bound.

\subsection{Local coarse gauge invariance}
\label{ca-gauge-section}

For a real coarse site field $\omega$ put $g=e^{i\omega}$,
$D=e^{iB}$ and
\begin{equation}\label{ca-coarse-gauge}
 D^g(c)=g(c_-)D(c)g(c_+)^{-1},\qquad
 B^g=\frac1i\log D^g.
\end{equation}
All logarithms use the common small branch. At coarse spacing
one the local analytic multiplication gives a constant
$C_c\ge1$ such that
\begin{equation}\label{ca-coarse-gauge-bound}
 \norm{B^g}_\infty\le C_c(b+w),\qquad
 w=\norm{\omega}_\infty.
\end{equation}
No additional regularity is required of $\omega$: coarse
first differences are bounded by $2w$.

Let $b_0,r_0$ be the ordinary-minimum radii of Notebook 15,
and write
$\norm{A_{\rm ord}(B)}_{N_\alpha}\le C_Ab$.
The smooth point-value extension there satisfies
\begin{equation}\label{ca-smooth-gauge-lift}
 (E_{\rm pt}\omega)|_{\rm coarse}=\omega,\qquad
 \norm{E_{\rm pt}\omega}_{C^{2,\alpha}}\le C_Ew.
\end{equation}
Set $h=e^{iE_{\rm pt}\omega}$, so $h|_{\rm coarse}=g$.
The fine gauge-action bound, valid for both $h$ and $h^{-1}$,
is
\begin{equation}\label{ca-fine-gauge-bound}
 \norm{A^h}_{N_\alpha}
       \le C_G(\norm A_{N_\alpha}+C_Ew).
\end{equation}

Choose one $\delta>0$ so that all logarithms and gauge actions
stay in their established domains and
\begin{equation}\label{ca-comparison-margins}
 C_c\delta<b_0/2,\qquad
 C_G(C_AC_c+C_E)\delta<r_0/2.
\end{equation}
For real $b+w<\delta$, ordinary gauge covariance makes
$U_{\rm ord}(B)^h$ feasible for coarse datum $B^g$.
Its potential norm is below $r_0$, so minimality at $B^g$
and exact gauge invariance give
\[
 V(B^g)\le\mathcal S_\eta(U_{\rm ord}(B)^h)=V(B).
\]
For the reverse comparison use
$U_{\rm ord}(B^g)^{h^{-1}}$. Its ordinary coarse average is
$(D^g)^{g^{-1}}=D$, and its norm is bounded by
\[
 C_G\{C_AC_c(b+w)+C_Ew\}<r_0.
\]
It is therefore an admissible competitor at $B$, proving
$V(B)\le V(B^g)$. We have proved
\begin{equation}\label{ca-gauge-invariance}
 \boxed{\qquad V(B^g)=V(B)\qquad(b+w<\delta).\qquad}
\end{equation}
Both comparisons lie strictly inside the same uniform minimum
radius. The reverse one uses the separately controlled minimizer
at $B^g$, not an assumed inverse of the first comparison map.

The two sides of \eqref{ca-gauge-invariance} are holomorphic
in $(B,\omega)$ on a uniform smaller complex product ball.
Their equality on an open set of the full real form implies
equality on that connected product ball by the finite-dimensional
holomorphic identity theorem, applied coordinate by coordinate.
For example, reduce $\delta$ to respect all analytic domains
and take $\norm B_\infty,\norm\omega_\infty<\delta/3$.
This extension is an analytic identity; it is not a minimum
claim for complex fields.

\subsection{The exact classical Ward identity}
\label{ca-ward-section}

Define the coarse left exponential differential and the
coarse covariant scalar difference by
\begin{equation}\label{ca-left-differential}
 \mathcal J_Bv=\int_0^1\operatorname{Ad}(e^{itB})v\,dt,
 \qquad
 (D_D\omega)(c)=\operatorname{Ad}_{D(c)}\omega(c_+)-\omega(c_-).
\end{equation}
The operator $\mathcal J_B$ is invertible on the fixed small
complex ball. The ordinary coarse-potential velocity under
$g(t)=e^{it\omega}$ is
\begin{equation}\label{ca-gauge-velocity}
 \left.\frac{d}{dt}B^{g(t)}\right|_{t=0}
             =-\mathcal J_B^{-1}D_D\omega.
\end{equation}
The coarse spacing is one, so there is no additional length
factor. Differentiate \eqref{ca-gauge-invariance} and use
\eqref{ca-first-derivative}:
\[
 0=-\langle m(B),\mathcal J_B^{-1}D_D\omega\rangle_1
   =-\langle D_D^T\mathcal J_B^{-T}m(B),\omega\rangle_1.
\]
The site field $\omega$ is arbitrary. Hence
\begin{equation}\label{ca-ward}
 \boxed{\qquad D_D^T\mathcal J_B^{-T}m(B)=0.\qquad}
\end{equation}
The real identity extends holomorphically. Thus the
left-coordinate classical coarse current is
$j_{\rm cl}(B)=\mathcal J_B^{-T}m(B)$, not generally $m(B)$.
No quantum determinant or quantum Ward assertion is involved.

\subsection{A local classical Hessian and an auxiliary positive inverse}
\label{ca-precision-section}

The weighted multiplier response \eqref{wr-dm-near-q}, together
with the envelope identity, gives
\begin{equation}\label{ca-hessian-weighted}
 Q_{\rm cl}(B):=D^2V(B)=Dm(B),\qquad
 \norm{[Q_{\rm cl}(B)-q]c}_{\infty,w}
       \le C_\kappa b\norm c_{\infty,w}
\end{equation}
for every admissible weight in \eqref{wr-weight}. As the
Hessian of a holomorphic scalar function, $Q_{\rm cl}(B)$
is symmetric for the physical bilinear transpose.
It differentiates in the ordinary coarse potential $B$.
A change to a based left-exponential fluctuation chart
requires the coordinate Jacobians and its background-current
second-derivative term; that differently based Hessian is
not being identified with $Q_{\rm cl}(B)$ here.

Use a fixed trace-orthonormal color and orientation basis in
each coarse cell. For a block kernel $T$ define the weighted
Schur norm
\begin{equation}\label{ca-schur}
 \norm T_\kappa=\max\left\{
   \sup_n\sum_m e^{\kappa d(n,m)}\norm{T_{nm}},\,
   \sup_m\sum_n e^{\kappa d(n,m)}\norm{T_{nm}}\right\}.
\end{equation}
Take $w_n(m)=e^{-\kappa d(n,m)}$ in
\eqref{ca-hessian-weighted}. In each scalar output row,
choose input phases to match its matrix entries. This bounds
its absolute row sum with the factor $e^{\kappa d(n,m)}$.
Summing the fixed number of component rows converts this to
matrix-block norms; transpose symmetry gives the column bound.
Consequently
\begin{equation}\label{ca-hessian-schur}
 \norm{Q_{\rm cl}(B)-q}_\kappa\le C'_{\kappa}b.
\end{equation}
The component factor depends only on dimension and color,
not on refinement or volume. These finite-regulator row tests
are bounded directions even for decaying weights.

Let $a$ lie in a fixed interval $0<a_-\le a\le a_+$.
The flat baseline estimates give, after fixing a sufficiently
small common $\kappa$,
\begin{equation}\label{ca-baseline-inverse}
 \norm{(q+aI)^{-1}}_\kappa\le r_\kappa
\end{equation}
uniformly. Set $P_{\rm cl}(B)=Q_{\rm cl}(B)+aI$.
For
\begin{equation}\label{ca-inverse-radius}
 C'_\kappa b\le a_-/2,\qquad
 r_\kappa C'_\kappa b\le1/2,
\end{equation}
the real positive baseline $q\ge0$ and the Schur operator
bound imply
\begin{equation}\label{ca-accretivity}
 \operatorname{Re}\langle x,P_{\rm cl}(B)x\rangle_{\rm Herm}
                 \ge\frac{a_-}{2}\norm x_2^2
\end{equation}
even at complex $B$. The numerical-range pairing here is
Hermitian; it is distinct from the analytic bilinear transpose.
The Schur norm is a matrix algebra norm, so the Neumann series
also gives
\begin{equation}\label{ca-inverse-bounds}
 \begin{split}
 \norm{P_{\rm cl}(B)^{-1}}_\kappa&\le2r_\kappa,\\
 \norm{P_{\rm cl}(B)^{-1}-(q+aI)^{-1}}_\kappa
                         &\le2r_\kappa^2C'_\kappa b.
 \end{split}
\end{equation}
This inverse is holomorphic and exponentially local in
coarse-cell distance, with uniform constants.

The positive $a$ is an auxiliary penalty added to a coarse
classical Hessian. This is not a physical Yang--Mills mass gap
or the precision of an integrated quantum effective action.
Spatial kernel decay does not assert strict dependence on
bounded background supports. Fluctuation determinants,
Jacobian factors, quantum interactions and their iteration
remain separate tasks. No finite total action is assigned
here to arbitrary bounded fields on the infinite lattice.

% No preamble. Finite-regulator quantum-domain interface; qg- labels.
\section{Size of the flat constrained Gaussian fluctuations}
\label{qg-section}

There is a quantitative reason why the preceding smooth-background
ball cannot itself be used as a uniform domain for integrating all
fine Gaussian fluctuations. We prove it already at the flat
reference, in four dimensions. Let $n_f=M^4n_c$ be the numbers
of fine and coarse sites of a compatible periodic regulator,
where $M=L^j\ge L\ge3$ and $\eta=M^{-1}$. Work in one fixed
real trace-orthonormal color component; the flat operator acts
identically on every color. Set
\[
 X=\ker Q\cap\ker(RD^T),\qquad K=d^Td.
\]
The positive restricted Maxwell form defines a centered real
Gaussian $Z\in X$ with density proportional to
$\exp[-\langle Z,KZ\rangle_\eta/2]$ relative to the
physical Euclidean measure on $X$. Its covariance operator
under the physical pairing is the constrained inverse
$\mathcal C$ used above.

\begin{lemma}\label{qg-dimension}
For each real color component,
\[
 \dim X=3(n_f-n_c),\qquad
 \operatorname{Tr}(d\mathcal C d^T)=\dim X.
\]
\end{lemma}
\begin{proof}
The map $Q$ on bonds is onto because $QH=I$, hence
$\dim\ker Q=4(n_f-n_c)$. The scalar projection $R$ has
rank $n_f-n_c$: its complementary projection has the full
coarse scalar rank $n_c$. The map $RD^T$ from $\ker Q$
onto $\operatorname{ran}R$ is surjective. Indeed the scalar
operator $T=G'R$ satisfies $Q_0T=0$ and $R\mathcal L T=R$,
where $\mathcal L=D^TD$. For $f\in\operatorname{ran}R$,
the bond field $DTf$ has $QDTf=D_cQ_0Tf=0$ and
$RD^TDTf=f$. Rank subtraction gives the dimension.

Let $J_X$ embed $X$ isometrically into the physical bond
Hilbert space, and let $K_X=J_X^TKJ_X$. Then
$\mathcal C=J_XK_X^{-1}J_X^T$. Cyclicity of the finite
operator trace gives
$\operatorname{Tr}(d\mathcal C d^T)
=\operatorname{Tr}(K_XK_X^{-1})=\dim X$.
\end{proof}

The physical covariance identity is
$\mathbb E\langle f,Z\rangle_\eta\langle g,Z\rangle_\eta
=\langle f,\mathcal Cg\rangle_\eta$.
Thus the covariance of the actual coordinate values is
$\eta^{-4}\mathcal C$, and that of the physical curls is
$\eta^{-4}d\mathcal C d^T$. There are $6n_f$ positive
plaquette components per color. Lemma~\ref{qg-dimension}
therefore gives the exact average
\begin{equation}\label{qg-curl-variance}
 \boxed{\quad
 \frac1{6n_f}\sum_p\mathbb E|(dZ)_p|^2
     =\frac{M^4-1}{2}.
 \quad}
\end{equation}
In particular at least one real curl coordinate has variance
at least $(M^4-1)/2$, irrespective of the periodic volume.
This conclusion uses the physical covariance normalization;
omitting $\eta^{-4}$ would give a different and incorrect scale.

Fix a deterministic background $A_0$, a radius $r>0$ and
a coupling factor $g>0$. The event
$N_\alpha(A_0+gZ)\le r$ bounds every physical first difference,
and hence bounds the selected curl coordinate by $C r$,
where the color-norm constant is fixed. It restricts a
one-dimensional centered Gaussian of standard deviation
at least $c\eta^{-2}$ to an interval of length at most
$C r/g$; the interval center may depend on $A_0$.
The maximum of the Gaussian density therefore proves
\begin{equation}\label{qg-small-ball}
 \boxed{\quad
 \mathbb P\{N_\alpha(A_0+gZ)\le r\}
       \le \min\left(1,\frac{C r\eta^2}{g}\right).
 \quad}
\end{equation}
The bound is uniform in the deterministic background,
refinement and compatible periodic volume. Thus for any
sequence $g=g_\eta$ with $g_\eta/\eta^2\to\infty$, a fixed
physical smooth ball has probability tending to zero under
this constrained flat Gaussian reference.

This is an obstruction to integrating all fluctuations inside
the minimizer's fixed $N_\alpha$ ball, not an obstruction to
a Yang--Mills theory. Equation~\eqref{qg-curl-variance} shows
that the dimensionless microscopic curl $\eta^2dZ$ has
bounded averaged variance. A quantum renormalization step
must instead control suitable rescaled fluctuation fields
and their tails while retaining smooth estimates for the
background. Averaged or component variance bounds alone
do not control a supremum over a growing number of variables,
nor do they bound interacting polymer activities. No
non-Gaussian probability estimate is inferred here.


\section{The next interface to quantum integration}
\label{nr16-next}

The classical background and its response are now controlled
on a uniform homogeneous small-field domain. The minimized
action has a true nonlinear local response and a coarse
Hessian close to the flat Maxwell precision. Adding an
auxiliary positive coarse penalty gives a controlled inverse,
but that penalty is not a physical mass and this inverse
is not yet the covariance produced by the complete quantum
renormalization step.

Exponential response alone does not give strict background
support for polymer activities. The physical Hessian,
gauge-coordinate factors and multiple-insertion local
decomposition require the separate constructions described
below. Regions outside the chart and preservation of uniform
complex-source bounds under quantum iteration remain essential.

Notebook 17 now proves a relative energy comparison for the actual
first-step reduced Hessian, retaining constraint curvature, and
joint Gaussian bad-block estimates on the fixed-step domain. Its
homogeneous flat current-scale theorem covers all previous averaging
depths. Notebook 18 handles overlapping local interaction supports
and strict background locality under a small local precision
decomposition. Notebook 19 supplies the actual physical
pivot Hessian and local measure factors. Notebook 20 closes
the bare Wilson small-field integration. Notebook 21 supplies
strict support for the minimized classical action, with
arbitrary fixed cardinality reserve, and extends one
classical-reference integration uniformly through earlier
refinement. The next-coarse support transport, inherited
quantum interactions and regional domains remain separate.

The continuum obligations remain those of Notebook 01:
a common nontrivial gauge-invariant observable hierarchy,
Euclidean covariance and positivity in the limit, appropriate
normalization and Yang--Mills identities, and the additional
vacuum and mass-gap properties. The classical theorems here
remove specific prerequisites; they do not establish those
quantum conclusions.

\begin{thebibliography}{9}
\bibitem{prior16}
\emph{Toward a four-dimensional Yang--Mills construction},
research Notebooks 8--15 in the present workspace,
16 October 2026. The weighted flat extension and inverse,
uniform nonlinear constraint chart, Wilson force, constrained
minimum and ordinary-average transfer are the prerequisites.
\bibitem{source16}
E.~Sussman and collaborators,
\href{https://github.com/ethanwsussman/Yang--Mills-LLM-experiment}{Yang--Mills-LLM-experiment},
deposited raw accounts I--III at commit
\texttt{e392ef45013a581664bab76b90867170b30fe10e}.
The provenance limitations of Notebook 15 remain in force.
\end{thebibliography}
\end{document}
